\chapter{The Bound}

Begin where the construction begins: with the first difference. Not a thing, not a number, not a space --- only
a \emph{this} that is not a \emph{that}. One distinction, held. Everything the instrument does, it does by iterating
that single act; and at every step it says out loud which kind of act it is performing --- posited, proved, read, or
reserved. It grants itself exactly one thing, the law of motion, and nothing else. That is the whole license. Watch
what follows from it, step by watched step.

The first move is the naming. To name is to draw a distinguishable --- to take a \emph{this} and separate it from
its \emph{that} by the cheapest decision that will settle them apart. Hand the construction the electron, and it does
not measure it; it \emph{names} it, and the name is a count. The electron reports its state for free: \(-1\), one
orientation of exactly three. This is not read off any dial. It falls out of the naming itself, because to be
distinguishable at all is already to occupy one of the three positions the loop admits. \textbf{This is a fact off the
build: the axiom footprint here is empty --- \texttt{\#print axioms} returns the empty list.} The state costs nothing
because the state \emph{is} the naming.

Now the residue. When the construction counts around the loop and returns to where it began, the count must close
--- the walk that leaves must be the walk that arrives, or the \emph{this} was never distinguishable from its
\emph{that} to begin with. Closure is the demand that \(1 = 0.999\ldots\): the tail of nines is not an approximation
creeping toward one, it is one, written the long way. The residue closes. And because it closes as a \emph{conserved
holonomy} and not as a quantity awaiting a unit, \(\alpha\) is dimensionless --- not a number that happens to have no
dimension, but a \emph{ratio the loop is forbidden to leave open}. Again: derived, empty-footprint, a fact off the
build.

Here is the deepest step, and I mark it plainly as \textbf{interpretation}, not as a proved theorem. Ask the electron
its number and it answers for free --- that was the state. But ask it to \emph{couple}, to name not another but
\emph{itself}, and the self-embrace is not free. \(\alpha\), on this reading, is the electron's \textbf{self-energy}:
the Richardson residual of the electron naming itself, the price of the fixed point. The state is what you are; the
coupling is what you pay to \emph{be} it. The construction can exhibit this residual as a structural necessity ---
that the self-naming leaves a remainder --- but the claim that this remainder \emph{is} the physical self-energy, the
radiative dressing, is a bridge I am building over the code, not a plank the code lays down.

From the closed residue the three holonomy states follow of themselves, and count-to-three seals them:
\(-1 / 0 / +1\) --- electron, the null identity, positron. The loop resolves to exactly these three and admits no
fourth, because a fourth would be a distinguishable the closure cannot carry. That count --- three and no more --- is
derived and empty-footprint. It is the entire \emph{shape} of the coupling, built with nothing borrowed.

And now the boundary, which is the point of the whole exercise. The construction can derive the structure --- state,
closure, dimensionlessness, the three, the self-embrace --- but it \textbf{cannot} produce the magnitude, the
\(137.036\) a Penning trap reads off the electron in any lab. It does not fail to. It \textbf{proves it cannot}, and
here is the logic, marked as interpretation where it leans on geometry: a machine that \emph{counts} measures the
\emph{trace} --- the Ricci part, volume, how a bundle of paths converges. That trace is exactly the structure it
derived. But the magnitude --- the size of the self-energy, the radiative dressing --- lives in the trace-\emph{free}
part, the Weyl curvature, pure shape. And a counter is blind to shape \emph{by construction}, because shape is
precisely the deformation that leaves the count unchanged. The electron pays its state in a currency the counter reads
and charges its coupling in a currency the counter cannot. Not hidden --- invisible, provably.

You can watch the cut get drawn to scale. The bisection constructs \(\alpha\) as Dedekind constructs a real: brackets
placed one bit at a time, the nines of \(0.999\ldots\) written one per step --- and it \textbf{halts at the floor}.
\textbf{That floor is a fact off the build: the bound is 60 bits.} The nines it places are the structure it sees; the
place it stops is the bound; the tail it cannot write is the magnitude. Near side, derivable step by watched step. Far
side, reserved.

That is the result, and it is heavier than the number. The construction did not come up short. It measured the
boundary of its own sight --- the whole structure on the near side of the floor, the magnitude provably on the far
side, and the cut between them drawn bit by bit. A statement about measurement itself. The reading is done.
